\documentclass[10pt]{article}

\usepackage{arxiv}
\usepackage[utf8]{inputenc}
\usepackage[T1]{fontenc}
\usepackage{hyperref}
\usepackage{url}
\usepackage{booktabs}
\usepackage{amsfonts}
\usepackage{amsmath}
\usepackage{graphicx}
\usepackage{enumitem}
\usepackage{listings}
\usepackage{xcolor}
\usepackage{longtable}
\usepackage{array}
\usepackage{multirow}
\usepackage{tabularx}
\usepackage{float}
\usepackage{tikz}

% Remove all colored frames from hyperlinks
\hypersetup{
  colorlinks=false,
  linkcolor=black,
  filecolor=black,
  citecolor=black,
  urlcolor=black,
  pdfborder={0 0 0}
}

% Code listing style
\lstdefinestyle{pythonstyle}{
  language=Python,
  basicstyle=\ttfamily\scriptsize,
  keywordstyle=\bfseries,
  commentstyle=\itshape,
  stringstyle=\ttfamily,
  frame=single,
  framerule=0.4pt,
  rulecolor=\color{black!30},
  backgroundcolor=\color{black!2},
  breaklines=true,
  breakatwhitespace=true,
  numbers=left,
  numberstyle=\tiny,
  numbersep=5pt,
  xleftmargin=15pt,
  showstringspaces=false,
  columns=flexible
}
\lstset{style=pythonstyle}

% Compact lists
\setlist{nosep, leftmargin=*}
\setlist[itemize]{label=\textbullet}
\setlist[enumerate]{label=\arabic*.}

\title{Federated Learning Physical AI Oncology Trials Unification}

\author{
  Kevin Kawchak \\
  CEO, ChemicalQDevice \\
  \href{https://orcid.org/0009-0007-5457-8667}{\textcolor{green!60!black}{\textbf{ORCID}}} \href{https://orcid.org/0009-0007-5457-8667}{https://orcid.org/0009-0007-5457-8667} \\
  \texttt{kevink@chemicalqdevice.com}
}

\date{February 27, 2026}

\begin{document}
\maketitle

\begin{abstract}
The transition from conventional software-only artificial intelligence to physical AI systems incorporating robotic hardware in oncology clinical trials represents a paradigm shift requiring unified infrastructure for privacy, regulation, cross-framework interoperability, and multi-organization cooperation. This paper presents the PAI Oncology Trial FL platform (v1.1.0), a comprehensive federated learning framework comprising 235 Python modules ($\sim$86,800 lines of code) that unifies five critical infrastructure pillars: (1) Privacy Infrastructure implementing all 18 HIPAA Safe Harbor identifiers with HMAC-SHA256 pseudonymization, (2) Regulatory Infrastructure spanning FDA, IRB, ICH-GCP, and multi-jurisdiction compliance across v0.6.0 and v0.9.1, (3) Cross-Framework Unification bridging NVIDIA Isaac Sim, MuJoCo, Gazebo, and PyBullet simulation environments, (4) Standards \& Benchmarking for Q1 2026 objectives including model conversion and registry pipelines, and (5) Multi-Organization Cooperation enabling federated training across academic medical centers, community hospitals, and pharmaceutical companies. End-to-end workflow demonstrations are presented across 31 example scripts, 6 agentic AI production examples implementing Model Context Protocol (MCP), ReAct reasoning, real-time monitoring, autonomous orchestration, safety-constrained execution, and RAG-based compliance. A triple AI peer review process (v0.9.4--v0.9.9) using sequential Codex-to-Claude Code review-fix cycles resolved 31/31 code recommendations at 100\% completion, establishing a dual-manufacturer trust benchmark for AI-generated clinical trial software. The platform demonstrates that unified federated learning infrastructure is a necessary precondition for transitioning the oncology industry to using robots in physical AI clinical trials.

\keywords{Federated Learning \and Physical AI \and Oncology Clinical Trials \and Digital Twins \and Regulatory Compliance \and HIPAA \and Cross-Framework Unification \and Agentic AI \and Model Context Protocol \and Differential Privacy}
\end{abstract}

\vspace{-0.5em}
\tableofcontents
\vspace{0.5em}

% =================================================================
\section{Introduction}
\label{sec:introduction}
% =================================================================

\subsection{Current AI in Federated Learning Oncology}

Federated learning (FL) has emerged as the dominant paradigm for privacy-preserving collaborative machine learning in oncology, enabling multiple institutions to jointly train AI models without exchanging raw patient data~\cite{mcmahan2017fedavg, rieke2020fl_medicine}. Leading cancer centers including Dana-Farber have launched federated AI platforms where models travel to each center's secure data environment~\cite{danafarberfederated2025}, and the European DigiONE consortium has demonstrated privacy-preserving learning health systems in precision oncology~\cite{digione2024}.

Current FL approaches in oncology typically address isolated aspects of the clinical trial workflow:
\begin{itemize}
  \item \textbf{Model training:} FedAvg~\cite{mcmahan2017fedavg}, FedProx~\cite{li2020fedprox}, and SCAFFOLD~\cite{karimireddy2020scaffold} for heterogeneous hospital data
  \item \textbf{Privacy:} Differential privacy with Gaussian noise mechanisms~\cite{dwork2014dp} and secure aggregation protocols
  \item \textbf{Brain tumor segmentation:} Multi-institutional FL for glioblastoma without patient data sharing~\cite{sheller2020fl_brain}
\end{itemize}

However, no existing open-source platform unifies federated learning, physical AI robotics, digital twins, regulatory compliance, privacy infrastructure, and clinical analytics into a single coherent framework for oncology clinical trials~\cite{kawchak2026fl}.

\subsection{Benefits of Upgrading to AI Robots in Oncology Trials}

The integration of physical AI---embodied robotic systems with onboard intelligence---into oncology clinical trials offers transformative benefits:
\begin{enumerate}
  \item \textbf{Precision in tumor biopsies and drug delivery:} Surgical robots executing procedures with sub-millimeter accuracy guided by digital twin predictions, as modeled in the PAI platform's \texttt{physical\_ai/} modules (v0.1.0)
  \item \textbf{Consistent monitoring:} Autonomous robotic systems performing 24/7 patient monitoring with real-time anomaly detection via multi-modal sensor fusion (v0.3.0)
  \item \textbf{Cross-site standardization:} Unified robot programming across trial sites eliminates procedural variability, enabled by the cross-framework unification layer (v0.3.0)
  \item \textbf{Federated model improvement:} Robotic telemetry from multiple sites feeds back into federated model training without raw data exchange, as demonstrated in \texttt{01\_mcp\_oncology\_server.py} (v0.5.0)
\end{enumerate}

\subsection{Projected Physical AI Benefits from Quantitative Codebase Data}

The PAI Oncology Trial FL platform (v1.1.0) provides quantitative evidence for physical AI readiness:

\begin{table}[H]
\centering
\caption{Platform quantitative metrics supporting physical AI oncology trial readiness}
\label{tab:quantitative}
\small
\begin{tabular}{@{}lrl@{}}
\toprule
\textbf{Metric} & \textbf{Value} & \textbf{Version} \\
\midrule
Total Python modules & 235 & v1.0.1 \\
Lines of code & $\sim$86,800 & v1.0.1 \\
Test files & 82 & v0.8.0+ \\
Test directories & 15 & v0.8.0+ \\
Example scripts & 31 & v0.5.0 \\
Agentic AI examples & 6 & v0.5.0 \\
CLI tools & 5 & v0.4.0 \\
Security audit findings resolved & 61 & v0.7.0 \\
Peer review recommendations resolved & 31/31 (100\%) & v0.9.4--v0.9.9 \\
Regulatory frameworks supported & 9 & v0.9.1 \\
HIPAA Safe Harbor identifiers & 18/18 & v0.6.0 \\
Simulation frameworks bridged & 4 & v0.3.0 \\
Tumor growth models & 4 & v0.4.0 \\
Supported Python versions & 3.10, 3.11, 3.12 & v0.3.0+ \\
\bottomrule
\end{tabular}
\end{table}

The upstream Physical AI Oncology Trials repository~\cite{kawchak2026physical} introduced the Unification Standard Level (USL) scoring framework evaluating 9 robotic systems across humanoid robots, surgical robots, and collaborative robots, with scores ranging from 3.4 (UFACTORY xArm7) to 7.4 (Franka Panda)~\cite{kawchak2026usl}. The PAI Oncology Trial FL platform extends this foundation with federated learning capabilities necessary for multi-site robot deployment.

% =================================================================
\section{Methods}
\label{sec:methods}
% =================================================================

\subsection{AI Development Tools and Models}

The PAI Oncology Trial FL platform was developed using the following AI systems:
\begin{itemize}
  \item \textbf{Primary development:} Anthropic Claude Code powered by Claude Opus 4.6~\cite{anthropic2026claude}---responsible for all code generation, documentation, and fix implementation across v0.1.0 through v1.1.0
  \item \textbf{Independent code review:} OpenAI ChatGPT-5.2-Codex~\cite{openai2026codex}---performed three sequential senior engineering review cycles (v0.9.4, v0.9.6, v0.9.8) generating 31 total recommendations
  \item \textbf{Paper generation:} This paper was primarily generated by Claude Code (Opus 4.6) based on Mr.\ Kawchak's hand-written prompt, which took approximately 2 hours to compose
\end{itemize}

\subsection{Development Timeline and Version History}

The codebase was developed over the period February 17--27, 2026, with the repository converted from private to public on February 27, 2026. Table~\ref{tab:timeline} summarizes the complete version history.

\begin{table}[H]
\centering
\caption{Complete release timeline with AI contributor roles}
\label{tab:timeline}
\small
\begin{tabular}{@{}llll@{}}
\toprule
\textbf{Version} & \textbf{Date} & \textbf{AI Role} & \textbf{Key Deliverable} \\
\midrule
v0.1.0 & 2026-02-17 & Claude Code & Initial FL platform \\
v0.2.0 & 2026-02-17 & Claude Code & FedProx, SCAFFOLD, data harmonization \\
v0.3.0 & 2026-02-17 & Claude Code & Unification framework, CI/CD \\
v0.4.0 & 2026-02-17 & Claude Code & Digital twins, 5 CLI tools \\
v0.5.0 & 2026-02-17 & Claude Code & 31 example scripts \\
v0.6.0 & 2026-02-18 & Claude Code & Privacy \& regulatory (9 modules) \\
v0.7.0 & 2026-02-18 & Claude Code & Security audit (61 findings) \\
v0.8.0 & 2026-02-18 & Claude Code & Test suite (82 files, 15 dirs) \\
v0.9.0 & 2026-02-18 & Claude Code & Clinical analytics (7 modules) \\
v0.9.1 & 2026-02-18 & Claude Code & Regulatory submissions (6 modules) \\
v0.9.2 & 2026-02-18 & Claude Code & Release documentation \\
v0.9.3 & 2026-02-19 & Claude Code & Visualization (30 charts) \\
v0.9.4 & 2026-02-19 & Codex & Peer review \#1: 12 recommendations \\
v0.9.5 & 2026-02-19 & Claude Code & Fix cycle \#1: 12/12 resolved \\
v0.9.6 & 2026-02-19 & Codex & Peer review \#2: 9 recommendations \\
v0.9.7 & 2026-02-19 & Claude Code & Fix cycle \#2: 9/9 resolved \\
v0.9.8 & 2026-02-19 & Codex & Peer review \#3: 10 recommendations \\
v0.9.9 & 2026-02-19 & Claude Code & Fix cycle \#3: 10/10 resolved \\
v1.0.0 & 2026-02-19 & Claude Code & Stable release \\
v1.0.1 & 2026-02-19 & Claude Code & Documentation update \\
v1.1.0 & 2026-02-27 & Claude Code & This paper and archival \\
\bottomrule
\end{tabular}
\end{table}

\textbf{Development statistics:}
\begin{itemize}
  \item \textbf{Total span:} 11 days (Feb 17--27, 2026)
  \item \textbf{Days with AI code updates:} 4 days (Feb 17, 18, 19, 27)
  \item \textbf{Days of AI code recommendations (Codex):} 1 day (Feb 19---v0.9.4, v0.9.6, v0.9.8)
  \item \textbf{Days of AI code fixes (Claude Code):} 4 days (Feb 17, 18, 19, 27---v0.1.0 through v1.1.0)
  \item \textbf{Human prompt authoring (Kevin Kawchak):} 18 prompts across development lifecycle, including a 2-hour hand-written prompt for this paper
  \item \textbf{Time from first prompt to PR submission:} Approximately 30 minutes (prompt submission: 2026-02-27 03:15 UTC, PR creation: 2026-02-27 03:45 UTC est.)
\end{itemize}

\subsection{Prompt-Driven Development Methodology}

All 18 development prompts are recorded in \texttt{prompts.md}~\cite{kawchak2026fl}. The methodology follows a single-prompt delivery pattern where each prompt generates a complete, self-contained release. Key prompts include:
\begin{itemize}
  \item \textbf{Prompt 1 (v0.1.0):} ``Integrate physical AI with federated learning platform instructions into \texttt{pai-oncology-trial-fl}. Utilize \texttt{physical-ai-oncology-trials} for physical AI implementation.''
  \item \textbf{Prompt 8 (v0.8.0):} ``Build comprehensive pytest-based test suite achieving full coverage of all Python modules.''
  \item \textbf{Prompt 13a (v0.9.4):} ``Provide comprehensive code reviews to find errors across the entire codebase'' (Codex)
  \item \textbf{Prompt 16 (v1.0.0):} ``Prove and push a v1.0.0 stable release by analyzing the codebase end-to-end''
\end{itemize}

